\chapter{Mean values and model checks for graphical flow}\label{sec:examples-periodic-graph-flow}
These examples test the conclusions of Sections \ref{sec:graphical-weighted-curvature} and \ref{sec:periodic-graph-convergence}. They use the graphical equation studied in \cite{EH89}, Section 3, but the periodic mean-value computations below are direct calculations.

\begin{example}[Curve graphs preserve their ordinary spatial average]\label{ex:curve-graph-mean-conservation}
Let $u(x,t)$ be a smooth periodic curve graph with spatial period $\ell$. Then
\[
 u_t=\frac{u_{xx}}{1+u_x^2}=\partial_x\arctan u_x,
 \qquad \frac d{dt}\int_0^\ell u(x,t)\,dx=0.
\]
If the solution is smooth for all positive time, its limiting height is exactly $\ell^{-1}\int_0^\ell u(x,0)dx$.
\end{example}
\begin{proof}
In dimension one Proposition \ref{prop:graphical-mcf} has the displayed coefficient. Differentiating $\arctan u_x$ gives the divergence expression. Integration over the period has no endpoint term because $u_x$ is periodic. If the flow is all-time smooth, Theorem \ref{thm:periodic-graph-conditional-convergence} gives uniform convergence to a constant, so its conserved integral identifies that constant. This conservation law comes from the one-dimensional equation; it is not a consequence merely of the area being translation invariant in height.
\end{proof}

\begin{example}[A two-dimensional graph changes its spatial average]\label{ex:surface-graph-changing-mean}
On the square torus with side length $2\pi$, set
\[
 f(x,y)=\cos x+\cos y+\alpha\cos(x+y),\qquad
 u(x,y,0)=\delta f(x,y),
\]
where $\alpha\ne0$ is fixed. For any smooth graphical solution with this initial value, its normalized average $\bar u(t)$ satisfies
\[
 \bar u'(0)=-\frac\alpha2\delta^3+O(\delta^5)
 \quad\text{as }\delta\to0.
\]
Thus it is nonzero for sufficiently small nonzero $\delta$. The construction also gives examples in every dimension at least two by making the height independent of the extra variables.
\end{example}
\begin{proof}
All averages in this proof are normalized by $(2\pi)^{-2}$. At the initial time the equation gives
\[
 \bar u'(0)=-\delta^3\operatorname{avg}
 \frac{f_i f_j f_{ij}}{1+\delta^2|Df|^2},
\]
with summation over $i,j\in\{1,2\}$, since the average of $\Delta f$ is zero. The bounded trigonometric derivatives justify replacing the denominator by one with an $O(\delta^2)$ error inside this average.

Periodic integration by parts gives
$\operatorname{avg}(f_i f_j f_{ij})=-\tfrac12\operatorname{avg}(|Df|^2\Delta f)$.
Writing $s=x+y$, one has
\[
 |Df|^2=\sin^2x+\sin^2y+2\alpha^2\sin^2s
       +2\alpha\sin s(\sin x+\sin y),
 \quad \Delta f=-\cos x-\cos y-2\alpha\cos s.
\]
Every product has zero average except
$-2\alpha\sin s\sin x\cos y$ and
$-2\alpha\sin s\sin y\cos x$.
Each has average $-\alpha/2$, using $\sin(x+y)=\sin x\cos y+\cos x\sin y$ and the averages of squared sines and cosines. Therefore $\operatorname{avg}(|Df|^2\Delta f)=-\alpha$ and the contracted Hessian average is $\alpha/2$. This gives the formula and its nonvanishing conclusion. In higher dimensions the zero extra derivatives leave the same normalized calculation. This computes the prescribed initial velocity of any smooth solution; existence itself is a separate parabolic input.
\end{proof}

\begin{example}[A single Fourier mode is not an exact nonlinear flow]\label{ex:graphical-single-mode-not-solution}
On a circle of length $2\pi$, the ansatz $u(x,t)=a(t)\cos x$ satisfies graphical mean curvature flow on a time interval only if $a$ vanishes identically. The linear heat ansatz $a(t)=a_0e^{-t}$ is an approximation at small amplitude, with exact residual
\[
 u_t-\frac{u_{xx}}{1+u_x^2}
   =-\frac{a(t)^3\cos x\sin^2x}{1+a(t)^2\sin^2x}.
\]
Its uniform residual is at most $|a(t)|^3$.
\end{example}
\begin{proof}
For the general ansatz the equation reads
$a'\cos x=-a\cos x/(1+a^2\sin^2x)$.
At $x=0$ it requires $a'=-a$. At any point where $\cos x\ne0$ and $\sin x\ne0$, the same equation, together with $a'=-a$, forces $a=0$. Thus at every time in the interval the amplitude vanishes. For $a'=-a$ without requiring an exact solution, subtraction gives the stated residual. Its bound follows from $|\cos x|,|\sin x|\leq1$ and the denominator being at least one. Exponential linear decay therefore does not furnish an exact nonlinear solution by itself.
\end{proof}
